\documentclass[12pt]{article}
\usepackage[utf8]{inputenc}
\usepackage[T1]{fontenc}
\usepackage{hyperref}
\title{Paleni Language Resources for K-12 Curriculum Development}
\author{Compiled by Claude Sonnet 5.5 (claude-sonnet-5-5) for Andr\'as Kornai}
\date{}
\begin{document}
\maketitle

\section{Introduction}
Paleni (ISO 639-3 pnl; also called Wara) is a small Niger-Congo language of Burkina Faso, spoken in the village of Faniagara. Wikipedia gives about 260 speakers (2012). It belongs to the Wara--Natyoro (Samu) group, tentatively placed in Gur by G\"uldemann (2018). Web search found very little educational material. No primers, textbooks or school programmes were found.

\section{Primary Educational Resources}
None found in search results.

\section{Linguistic and Reference Materials}
\begin{itemize}
\item Paleni language, Wikipedia: \url{https://en.wikipedia.org/wiki/Paleni_language}
\item Wara--Natyoro languages, Wikipedia: \url{https://en.wikipedia.org/wiki/Wara%E2%80%93Natyoro_languages}
\item University of Michigan Deep Blue data download that appeared in results for Paleni (reportedly documentation data by Jeffrey Heath; contents not verified): \url{https://deepblue.lib.umich.edu/data/downloads/p8418p57f}
\item ASJP word list entry: \url{https://asjp.clld.org/languages/PALEN}
\end{itemize}

\section{Religious and Community Resources}
\begin{itemize}
\item Joshua Project profile (translation status listed as started; not verified further): \url{https://joshuaproject.net/people_groups/21002}
\item SIL Burkina page that appeared in results (content not verified): \url{https://sil-burkina.org/node/246}
\end{itemize}

\section{Digital and Audio Resources}
No audio or software resources were found.

\section{Curriculum Development Context}
Material in Paleni is essentially limited to linguistic documentation. A curriculum would need to start from community-led orthography and materials development, drawing on the documentation data above.
\end{document}